\documentclass[conference]{IEEEtran}
\IEEEoverridecommandlockouts
\usepackage{cite}
\usepackage{amsmath,amssymb,amsfonts}
\usepackage{algorithmic}
\usepackage{graphicx}
\usepackage{textcomp}
\usepackage{xcolor}
\def\BibTeX{{\rm B\kern-.05em{\sc i\kern-.025em b}\kern-.08em
    T\kern-.1667em\lower.7ex\hbox{E}\kern-.125emX}}

\begin{document}

\title{Security and Software Application Project Report}

\author{\IEEEauthorblockN{Luca Sforza}
\IEEEauthorblockA{\textit{Sapienza Università di Roma} \\
sforza.2050030@studenti.uniroma1.it}
}

\maketitle

\begin{abstract}
Il presente elaborato illustra il lavoro svolto per il progetto d'esame del corso di Security and Software Application presso Sapienza Università di Roma.
L'obiettivo principale riguarda l'analisi e la correzione di vulnerabilità in smart contract Solidity forniti nelle specifiche di progetto.
Per lo sviluppo è stata adottata la metodologia Agile, suddividendo il lavoro in iterazioni incrementali al fine di garantire un controllo costante sulla qualità del codice e sulla sicurezza degli invarianti di sistema.
\end{abstract}

% Nota: In IEEEtran solitamente si usa l'inglese per Abstract e Keywords, 
% ma se il corso richiede l'italiano, il testo sopra è corretto.

\section{Introduction}

L'obiettivo di questo progetto consiste nell'analisi, nella correzione e nella messa in sicurezza degli smart contract forniti nella traccia d'esame.
Il lavoro si concentra sull'identificazione delle vulnerabilità presenti nel codice sorgente iniziale, verificate tramite la definizione e il controllo delle invarianti descritte nelle specifiche.

La fase di analisi è stata condotta seguendo un approccio \textit{Model Driven}, basato su UML (\textit{Unified Modeling Language}), un formalismo particolarmente adatto per la modellazione di sistemi complessi come quelli basati su Solidity.

\section{Refactoring e Pulizia del Codice}

Il codice fornito inizialmente presentava numerose criticità strutturali e stilistiche. 
In particolare, la visibilità degli attributi era impostata sul valore di default (spesso non intenzionale) e l'indentazione non seguiva alcuno standard di formattazione riconosciuto.

Per identificare e risolvere sistematicamente queste problematiche, sono stati utilizzati strumenti di analisi statica (linter).

\subsection{Ambiente di Sviluppo}

Per questo progetto è stato utilizzato il framework di sviluppo \textbf{Foundry}, che fornisce una suite di strumenti essenziali per la sicurezza e il testing di smart contract:

\begin{itemize}
    \item \textbf{Forge}: Utilizzato per compilare, testare e distribuire i contratti.
    \item \textbf{Anvil}: Un nodo locale Ethereum per il testing rapido.
    \item \textbf{Cast}: Strumento per interagire con il contratto tramite riga di comando.
\end{itemize}

\subsection{Analisi dei Problemi Logici}

Oltre alle problematiche sintattiche, il codice presentava vulnerabilità logiche significative. Queste sono state individuate grazie all'approccio \textit{Model Driven}.

L'utilizzo di diagrammi delle classi UML permette una traduzione naturale verso il codice Solidity. In letteratura sono proposti diversi meta-modelli per facilitare questa transizione.

% TODO: Inserire qui il riferimento alla figura: \begin{figure} ... \end{figure}

Il diagramma delle classi UML definisce i vincoli del sistema sotto forma di invarianti, le quali devono rimanere vere durante l'intera esecuzione delle transazioni.

Un esempio critico riscontrato riguarda la gestione dell'età del contribuente. Inizialmente, questa era memorizzata come un numero intero da incrementare manualmente ogni anno. Tale approccio è prono a errori e disallineamenti.
La soluzione adottata prevede invece l'utilizzo del timestamp del blocco in cui viene eseguita la transazione. In Solidity, è possibile accedere ai metadati della transazione (come il mittente \texttt{msg.sender}, il valore inviato \texttt{msg.value}, e il timestamp \texttt{block.timestamp}) tramite le variabili globali del linguaggio.

Pertanto, l'età non viene più salvata come stato mutabile, ma calcolata dinamicamente come metodo del contratto, garantendo coerenza temporale.

Per la verifica della correttezza del software, anziché procedere con una verifica formale matematica completa, si è optato per il \textbf{Fuzz Testing}.
Molti vincoli UML sono stati implementati direttamente tramite la sintassi di Solidity (modificatori di visibilità, tipi di dato, ereditarietà). I vincoli più complessi (invarianti inter-contratto) sono stati invece tradotti in metodi di verifica che analizzano lo stato globale del contratto per confermare il rispetto delle specifiche.

\section{Implementazione dei Vincoli: Esercizio 1}

Il primo esercizio richiede l'implementazione del vincolo di reciprocità coniugale ("spouse").
Formalmente, si deve verificare la seguente formula della logica del primo ordine:

\begin{equation}
\label{V_spouse}
    \forall a \in \text{Taxpayer}, \forall b \in \text{Person}: \text{spouse}(a,b) \implies \text{spouse}(b,a)
\end{equation}

In termini pratici: se un contribuente $a$ è registrato come sposo di $b$, allora $b$ deve necessariamente avere $a$ registrato come proprio sposo.

A differenza delle semplici invarianti di stato, questo vincolo implica una dipendenza circolare che non può essere modellata con la sola sintassi dichiarativa di Solidity. È necessario garantire che l'invariante \ref{V_spouse} sia preservata atomicamente durante l'esecuzione di qualsiasi metodo che modifichi lo stato coniugale (es. funzione \texttt{marry}).

Se il Taxpayer $A$ esegue il metodo \texttt{marry} verso $B$, il contratto deve forzare anche l'aggiornamento dello stato di $B$ verso $A$ all'interno della stessa transazione. Questo scenario genera una potenziale ricorsione infinita se $B$ tentasse di richiamare \texttt{marry} su $A$.
La soluzione implementata prevede l'introduzione di un caso base nella logica di ricorsione per terminare l'esecuzione una volta stabilita la relazione bidirezionale.

La correttezza di questa implementazione è stata validata utilizzando \textbf{Echidna} per il fuzzing, creando una proprietà (property) specifica che modella l'invariante \ref{V_spouse} e tenta di falsificarla con input casuali.

\section{Esercizio 2}

\begin{equation}
    \forall a,b ...
\end{equation}

\end{document}
